\documentclass[10pt,a4paper]{amsart}
\usepackage[margin=19mm]{geometry}
\usepackage{amsmath,amssymb,mathtools}
\usepackage[scr=rsfs,cal=euler]{mathalfa}
\usepackage{xfrac}
\usepackage[unicode,hidelinks,bookmarks=false]{hyperref}
\newcommand{\ie}{i.e.\ }
\newcommand{\cf}{cf.\ }
\newcommand{\resp}{respectively\ }
\newcommand{\Subsection}{\subsection}
\providecommand{\nobreakdash}{\nobreak\hbox{-}}
\setlength{\emergencystretch}{2em}
\multlinegap0pt
\usepackage[shortlabels]{enumitem}
\usepackage{mathtools}
\usepackage{amssymb}
\newtheorem{lem}{Lemma}
\newtheorem{claim}{Claim}
\newcommand{\NN}{\ensuremath{\mathbb N}}
\newcommand{\cB}{\mathcal{B}}
\newcommand{\floor}[1]{\mbox{$\left\lfloor #1 \right\rfloor$}}
\title{On the Thickness of Infinite Generalized Sidon Sets, I}
\author{Kevin O'Bryant}
\date{Source: arXiv:2606.28651v3 (CC BY 4.0)}
\begin{document}
\maketitle

\noindent\emph{Source and licence.} On the Thickness of Infinite Generalized Sidon Sets, I. Version arXiv:2606.28651v3 (CC BY 4.0), distributed by arXiv under the Creative Commons Attribution 4.0 International licence, http://creativecommons.org/licenses/by/4.0/, as recorded in the arXiv licence metadata for this exact version and shown on the abstract page https://arxiv.org/abs/2606.28651v3. Full licence notice: \texttt{licenses/arxiv-2606.28651-CC-BY-4.0.txt}.\par\smallskip

\noindent\textit{Editorial scope.} Counted unit: the article's lem lem:energy (source0.tex lines 132--151). The article's complete original proof of it is delivered with it (source0.tex lines 153--281, one \texttt{proof} environment of 129 lines). No mathematics is written, repaired or re-derived: the statement and the proof are the article's own lines, in the article's own notation, and no retained line is substituted or re-flowed.  No further source-local context is needed: every cross-reference of every retained line resolves inside the two delivered spans.  Nothing was omitted from any retained span, so the package carries no omission record.  No retained line contains a citation, so the wrapper carries no bibliography and no bibliography entry.  No retained line contains a figure environment, an image-inclusion command or drawing code, so no image is delivered and the package declares no asset dependency.  Editorial formatting only, in addition: an AMS \texttt{amsart} preamble that restates the article's own declarations and macros used by the retained lines, namely the environments \texttt{lem}, \texttt{claim} on separate counters, together with the standard packages those lines need.  The retained environments print in this document as Lemma 1, Claim 1 in order of appearance, whereas the article prints the same items with the numbering of its own full document.  The complete native source of the article as distributed in the arXiv e-print for this exact version is bundled unchanged as \texttt{sources/203866/source0.tex}.
\begin{lem}\label{lem:energy}
Let $\cB\subseteq\NN$ have counting function $B(n)$.
Let $M=M(N)$ satisfy, as $N\to\infty$,
\begin{enumerate}[label=\textit{(\roman*)},noitemsep]
  \item\label{eng:window} $B((M+1)N)=o(N)$;
  \item\label{eng:ratio} $\log M/\log N =1+o( 1 )$;
  \item\label{eng:head} $B(N)=o(\sqrt{N \log N})$.
\end{enumerate}
Suppose there is a constant $c$ such that there is an
offset $t^\ast=t^\ast(N)\in[0,N)$ whose block counts
  \begin{equation*}\label{hyp:blocks}
    F_\ell \coloneqq B(t^\ast+\ell N)- B(t^\ast+(\ell-1)N)
  \end{equation*}
satisfy (as $N\to\infty$)
  \begin{equation}\label{hyp:energy}
    \sum_{\ell=1}^{M}\binom{F_\ell}{2}\le cN+o(N).
  \end{equation}
Then
  \[\liminf_{m\to\infty}\frac{B(m)}{\sqrt{m/\log m}}\le \sqrt{\frac{8c}{\log 2}}.\]
\end{lem}

\begin{proof}
Assume $N\ge 3$, and set $\tau_N,\widetilde{\tau}_N$ as
  \begin{equation}\label{def tau}
    \tau_N \coloneqq \inf_{n\ge N} B(n)\sqrt{\frac{\log n}{n}},
    \qquad 
    \widetilde{\tau}_N \coloneqq \min\{\tau_N , 1+\sqrt{\frac{8c}{\log 2}}\},
  \end{equation}
so that $\tau_N$ is nondecreasing, $\lim_N\tau_N=\liminf_m \frac{B(m)}{\sqrt{m/\log m}}$, and $\widetilde{\tau}_N$ is bounded.
As $\tau_N$ is monotone, it suffices to show $\tau_N\le
\sqrt{8c/\log2}$ for arbitrarily large $N$. 
Given $N$, take $t^\ast$ as provided, and consider the ``energy''
  \[E\coloneqq \sum_{\ell=1}^{M}F_\ell^2
      =2\sum_{\ell=1}^{M}\binom{F_\ell}{2}+\sum_{\ell=1}^{M}F_\ell.\]

The hypothesis~\eqref{hyp:energy} bounds the first sum by $2cN+o(N)$, and the second sum telescopes:
  \[\sum_{\ell=1}^{M}F_\ell=B(t^\ast+\floor{M}N)-B(t^\ast)\le B((M+1)N) = o(N)\]
by~\ref{eng:window}. Hence 
\begin{equation}\label{eq:upper}
  E\le 2cN+o(N).
\end{equation}

The lower bound on $E$ uses Cauchy's Inequality with the weights
\[
   w_\ell\coloneqq \begin{cases}
    \left( \ell\,\log \ell N \right)^{-1/2} & 1\le\ell\le M; \\
    0 & \text{otherwise.}
   \end{cases}
\]
By Cauchy's inequality,
\begin{equation}\label{eq:CS}
   E \coloneqq \sum_{\ell=1}^{M}F_\ell^{2}
   \ge \frac{\bigl(\sum_{\ell=1}^{M}w_\ell F_\ell\bigr)^{2}}
           {\sum_{\ell=1}^{M}w_\ell^{2}} .
\end{equation}
We need an upper bound on the denominator ${\sum_{\ell=1}^{M}w_\ell^{2}}$, and a lower bound on the numerator $\bigl(\sum_{\ell=1}^{M}w_\ell F_\ell\bigr)^{2}$.

\begin{claim}\label{eq:den}
    \(\displaystyle {\sum_{\ell=1}^{M}w_\ell^{2}} \le \log 2 + o(1)\).
\end{claim}
\begin{proof}[Proof of Claim~\ref{eq:den}]
As $x\mapsto\bigl(x\,\log xN \bigr)^{-1}$ is decreasing,
\begin{align*}
   \sum_{\ell=1}^{M}w_\ell^{2}
   &=\sum_{\ell=1}^{M}\frac{1}{\ell\,\log \ell N}\\
   &\le\frac{1}{\log N}+\int_{1}^{M}\frac{dx}{x\,\log xN}\\
   &=\frac{1}{\log N}+\log\frac{\log MN}{\log N}.
\end{align*}
Now, by~\ref{eng:ratio}, the ratio $\log(MN)/\log(N)\to 2$, so that 
  \(   \sum_{\ell=1}^{M}w_\ell^{2}\le o(1)+\log 2.\)
\end{proof}

\begin{claim}\label{eq:num}
    \(\displaystyle \bigg(\sum_{\ell=1}^{M}w_\ell F_\ell\bigg)^{2} \ge \widetilde{\tau}_N^2\left(\frac{\log 2}{2}\right)^2 N+ o(N).\)
\end{claim}
\begin{proof}[Proof of Claim~\ref{eq:num}]
To ease the notation visually, set $\beta_i=B(T+iN)$, and note that $F_\ell = \beta_\ell - \beta_{\ell-1}$. Also, set $M'=\floor{M}$. Rearranging the summation gives
\begin{align*}
   \sum_{\ell=1}^{M'}w_\ell F_\ell
   &=\sum_{\ell=1}^{M'}w_\ell(\beta_\ell-\beta_{\ell-1})\\
   &=\sum_{\ell=1}^{M'-1}\beta_\ell\,(w_\ell-w_{\ell+1})+w_{M'} \beta_{M'}-w_1\beta_0 \\
   &\ge\sum_{\ell=1}^{M'-1}\beta_\ell\,(w_\ell-w_{\ell+1})-w_1\beta_0.
\end{align*}
Now, we have from~\ref{eng:head}
    \[w_1\beta_0=\frac{B(T)}{\sqrt{\log N}}\le \frac{B(N)}{\sqrt{\log n}} = \frac{o(\sqrt{N \log N})}{\sqrt{\log N}} =o(\sqrt N)\]
For $1\le\ell\le M'-1$ we have 
$T+\ell N\ge\ell N\ge N$, so the definitions of $\tau_N,\widetilde{\tau}_N$ on Line \eqref{def tau} yields
  \[\beta_\ell\ge\tau_N \sqrt{\ell N/\log \ell N}  \ge \widetilde{\tau}_N \sqrt{\ell N/\log \ell N} .\] 
Writing 
    \[v_\ell\coloneqq \sqrt{\ell N/\log \ell N}\] 
and using $w_\ell-w_{\ell+1}\ge0$,
\[
   \sum_{\ell=1}^{M'}w_\ell F_\ell
   \ge\widetilde{\tau}_N\sum_{\ell=1}^{M'-1}v_\ell(w_\ell-w_{\ell+1})-o(\sqrt N) .
\]
By rearranging the summation again, with $v_0\coloneqq 0$,
\[
   \sum_{\ell=1}^{M'-1}v_\ell(w_\ell-w_{\ell+1})
   =\sum_{\ell=1}^{M'-1}w_\ell(v_\ell-v_{\ell-1})-w_{{M'}}v_{M' -1}.
\]
Since the $v$ sequence is positive and strictly increasing and the $w$ sequence is positive, we have
    \begin{equation*}
    0 < w_{M'}v_{M'-1} 
    \le w_{M'}v_{M'}
    = \frac{\sqrt{N}}{\log M'N} 
    =o(\sqrt N).
    \end{equation*}

At this point, we have
    \[\sum_{\ell=1}^{M} w_\ell F_\ell \ge o(\sqrt N)+ \widetilde{\tau}_N \sum_{\ell=1}^{M'-1} w_\ell(v_\ell-v_{\ell-1}),\]
where $w_\ell$ and $v_\ell$ are explicit sequences with easily handled properties. Set $W(x)=(x \log xN)^{-1/2}$ and $V(x)=(xN/\log xN)^{1/2}$. For $N\ge e^2$, both $W$ and $V'$ are positive and decreasing on $x\ge 1$. 
For $2\le \ell\le M'-1$, since $V'$ is decreasing and $W(x)\le W(\ell)=w_\ell$ for $x\ge\ell$,
  \[w_\ell(v_\ell - v_{\ell-1}) = w_\ell\int_{\ell-1}^{\ell} V'(x)\,dx
    \ge w_\ell\int_{\ell}^{\ell+1} V'(x)\,dx
    \ge \int_{\ell}^{\ell+1} W(x)V'(x)\,dx.\]
For $\ell=1$ (recalling $v_0=0$), we have $V(2)\le\sqrt2\,V(1)\le 2V(1)$, so that
  \[w_1(v_1-v_0) = W(1)V(1) \ge W(1)\bigl(V(2)-V(1)\bigr) \ge \int_1^2 W(x)V'(x)\,dx.\]
Summing over $1\le\ell\le M'-1$,
\begin{align*}
    \sum_{\ell=1}^{M'-1} w_\ell(v_\ell-v_{\ell-1})
    &\ge \int_1^{M'} W(x) V'(x)\,dx \\
    &=\frac{\sqrt{N}}{2} \left(\log\left(1+\frac{\log M'}{\log N}\right)- \frac{\log M'}{\log(N) \log M'N}\right)
\end{align*}
By~\ref{eng:ratio}, we have $\log M' /\log N = 1+o(1)$ and so
\begin{align*}
  \sum_{\ell=1}^{M'} w_\ell(v_\ell-v_{\ell-1})
    &=\frac{\sqrt{N}}{2} \left(\log(2+o(1))-o(1)\right) \\
    &=\frac{\log 2}{2}\sqrt{N} + o(\sqrt N).
\end{align*}
Thus,
    \begin{align*}
    \left(\sum_{\ell=1}^{M} w_\ell F_\ell\right)^2
    &\ge \left( o(\sqrt{N})+ \widetilde{\tau}_N \left(\frac{\log 2}{2}\sqrt{N} + o(\sqrt N)\right)\right)^2 \\
    &= \widetilde{\tau}_N^2\frac{\log^2 2}{4} N + o(N).
    \end{align*}
Claim~\ref{eq:num} is verified.
\end{proof}

Using Claim~\ref{eq:den} and Claim~\ref{eq:num} in Cauchy's Inequality~\eqref{eq:CS}, we get a lower bound on the energy:
\begin{equation}\label{eq:lower}
   E\ge\frac{\widetilde{\tau}_N^{2}\bigl(\tfrac{\log2}{2}\bigr)^{2}N\,(1+o(1))}
            {(\log2)(1+o(1))}
   =\widetilde{\tau}_N^{2}\,\frac{\log2}{4}\,N+o(N).
\end{equation}
Combine the upper bound on Line~\eqref{eq:upper} and lower bound on Line~\eqref{eq:lower}, take $N\to\infty$, and we have
    \[ \left(\lim_{N\to\infty} \widetilde{\tau}_N\right)^2\frac{\log 2}{4} \le 2c,\]
where the limit exists because $\widetilde{\tau}_N$ is nondecreasing and bounded.
This implies that $\widetilde{\tau}_N\le \sqrt{\frac{8c}{\log 2}}$, so that by definition $\widetilde{\tau}_N=\tau_N$.
The inequality $\tau_N \le \sqrt{\frac{8c}{\log 2}}$ is the conclusion of Lemma~\ref{lem:energy}.
\end{proof}

\end{document}
